% Archivo compartido: no compilar directamente.
% Use clase_3_estudiante.tex o clase_3_docente.tex como archivo raíz.
\ifdefined\documentoSemanal
\else
\usepackage{tikz}
\usepackage{estilo_apuntes}
\usetikzlibrary{babel}
\encabezado{Semana 1 · Clase 3}{Material de clase}

\begin{document}
\fi

\ifdefined\documentoSemanal
\else
\tituloclase{Composición y comportamiento de funciones}{Transformaciones, composición, simetría y monotonía}{Semana 1 · Clase 3 · 2 horas}

\begin{objetivos}
\begin{itemize}
  \item Describir transformaciones elementales de gráficos.
  \item Operar y componer funciones considerando sus dominios.
  \item Reconocer funciones pares e impares.
  \item Interpretar crecimiento, decrecimiento y monotonía de una función.
\end{itemize}
\end{objetivos}

\begin{resumen}
Las transformaciones producen nuevos gráficos a partir de funciones conocidas. Las operaciones y composiciones relacionan reglas, pero obligan a controlar los dominios. La paridad y la monotonía describen propiedades globales o por intervalos.
\end{resumen}
\fi

\section*{Transformaciones de gráficos}

\begin{proposicion}[de transformación]
Sea $f:A\to\mathbb{R}$ y sean $h,k\in\mathbb{R}$ y $a,b\in\mathbb{R}\setminus\{0\}$. Para comparar los seis casos se utilizará como función base
\[
f:[0,2]\to[0,\sqrt2],\qquad f(x)=\sqrt{x}.
\]
En cada gráfica, el trazo gris discontinuo representa $f$ y el trazo azul representa la función transformada $g$.

\begin{enumerate}
\item \textbf{Traslación vertical.} La fórmula general es
\[
g(x)=f(x)+k.
\]
El gráfico se traslada $k$ unidades hacia arriba si $k>0$ y $|k|$ unidades hacia abajo si $k<0$. Para $k=1$:
\begin{center}
% Traslación vertical.
\begin{tikzpicture}[x=0.72cm,y=0.48cm,every node/.style={font=\scriptsize}]
  \draw[Gris!70] (-0.25,0)--(2.55,0) node[right] {$x$};
  \draw[Gris!70] (0,-0.15)--(0,2.75) node[above] {$y$};
  \draw[Gris,dashed,thick,domain=0:2,samples=40] plot (\x,{sqrt(\x)});
  \draw[Azul,very thick,domain=0:2,samples=40] plot (\x,{sqrt(\x)+1});
\end{tikzpicture}
\[
g(x)=f(x)+1,\qquad
\operatorname{Dom}(g)=[0,2],\qquad
\operatorname{Rec}(g)=[1,1+\sqrt2].
\]
\end{center}

\item \textbf{Traslación horizontal.} La fórmula general es
\[
g(x)=f(x-h).
\]
El gráfico se traslada $h$ unidades hacia la derecha si $h>0$ y $|h|$ unidades hacia la izquierda si $h<0$. Para $h=1$:
\begin{center}
% Traslación horizontal.
\begin{tikzpicture}[x=0.72cm,y=0.48cm,every node/.style={font=\scriptsize}]
  \draw[Gris!70] (-0.25,0)--(3.45,0) node[right] {$x$};
  \draw[Gris!70] (0,-0.15)--(0,1.8) node[above] {$y$};
  \draw[Gris,dashed,thick,domain=0:2,samples=40] plot (\x,{sqrt(\x)});
  \draw[Azul,very thick,domain=1:3,samples=40] plot (\x,{sqrt(\x-1)});
\end{tikzpicture}
\[
g(x)=f(x-1),\qquad
\operatorname{Dom}(g)=[1,3],\qquad
\operatorname{Rec}(g)=[0,\sqrt2].
\]
\end{center}

\item \textbf{Reflexión respecto del eje $x$.} La fórmula es
\[
g(x)=-f(x).
\]
Cada punto $(x,y)$ del gráfico de $f$ se transforma en $(x,-y)$:
\begin{center}
% Reflexión respecto del eje x.
\begin{tikzpicture}[x=0.72cm,y=0.48cm,every node/.style={font=\scriptsize}]
  \draw[Gris!70] (-0.25,0)--(2.55,0) node[right] {$x$};
  \draw[Gris!70] (0,-1.75)--(0,1.75) node[above] {$y$};
  \draw[Gris,dashed,thick,domain=0:2,samples=40] plot (\x,{sqrt(\x)});
  \draw[Azul,very thick,domain=0:2,samples=40] plot (\x,{-sqrt(\x)});
\end{tikzpicture}
\[
g(x)=-f(x),\qquad
\operatorname{Dom}(g)=[0,2],\qquad
\operatorname{Rec}(g)=[-\sqrt2,0].
\]
\end{center}

\item \textbf{Reflexión respecto del eje $y$.} La fórmula es
\[
g(x)=f(-x).
\]
Cada punto $(x,y)$ del gráfico de $f$ se transforma en $(-x,y)$:
\begin{center}
% Reflexión respecto del eje y.
\begin{tikzpicture}[x=0.72cm,y=0.48cm,every node/.style={font=\scriptsize}]
  \draw[Gris!70] (-2.35,0)--(2.35,0) node[right] {$x$};
  \draw[Gris!70] (0,-0.15)--(0,1.75) node[above] {$y$};
  \draw[Gris,dashed,thick,domain=0:2,samples=40] plot (\x,{sqrt(\x)});
  \draw[Azul,very thick,domain=-2:0,samples=40] plot (\x,{sqrt(-\x)});
\end{tikzpicture}
\[
g(x)=f(-x),\qquad
\operatorname{Dom}(g)=[-2,0],\qquad
\operatorname{Rec}(g)=[0,\sqrt2].
\]
\end{center}

\item \textbf{Escala vertical.} La fórmula general es
\[
g(x)=a f(x),\qquad a\neq0.
\]
Las ordenadas se multiplican por $|a|$; si $a<0$, también hay una reflexión respecto del eje $x$. Para $a=2$:
\begin{center}
% Escala vertical.
\begin{tikzpicture}[x=0.72cm,y=0.40cm,every node/.style={font=\scriptsize}]
  \draw[Gris!70] (-0.25,0)--(2.55,0) node[right] {$x$};
  \draw[Gris!70] (0,-0.15)--(0,3.1) node[above] {$y$};
  \draw[Gris,dashed,thick,domain=0:2,samples=40] plot (\x,{sqrt(\x)});
  \draw[Azul,very thick,domain=0:2,samples=40] plot (\x,{2*sqrt(\x)});
\end{tikzpicture}
\[
g(x)=2f(x),\qquad
\operatorname{Dom}(g)=[0,2],\qquad
\operatorname{Rec}(g)=[0,2\sqrt2].
\]
\end{center}

\item \textbf{Escala horizontal.} La fórmula general es
\[
g(x)=f(bx),\qquad b\neq0.
\]
Las abscisas se multiplican por $1/|b|$; si $b<0$, también hay una reflexión respecto del eje $y$. Para $b=2$:
\begin{center}
% Escala horizontal.
\begin{tikzpicture}[x=0.72cm,y=0.48cm,every node/.style={font=\scriptsize}]
  \draw[Gris!70] (-0.25,0)--(2.55,0) node[right] {$x$};
  \draw[Gris!70] (0,-0.15)--(0,1.75) node[above] {$y$};
  \draw[Gris,dashed,thick,domain=0:2,samples=40] plot (\x,{sqrt(\x)});
  \draw[Azul,very thick,domain=0:1,samples=40] plot (\x,{sqrt(2*\x)});
\end{tikzpicture}
\[
g(x)=f(2x),\qquad
\operatorname{Dom}(g)=[0,1],\qquad
\operatorname{Rec}(g)=[0,\sqrt2].
\]
\end{center}
\end{enumerate}

Las transformaciones verticales conservan el dominio $A$. En cambio,
\[
\operatorname{Dom}(f(x-h))=\{x:x-h\in A\},\qquad
\operatorname{Dom}(f(bx))=\{x:bx\in A\}.
\]
Por tanto, una transformación del argumento también transforma el dominio.
\end{proposicion}

\begin{ejemplo}[transformaciones]
Determine el mayor conjunto $D\subseteq\mathbb{R}$ para que la relación $g(x)=-2\sqrt{x-3}+1$ defina una función real $g:D\to\mathbb{R}$. Describa su gráfico y determine su imagen.
\end{ejemplo}
\ifmostrarSoluciones
\begin{solucion}
Desde $y=\sqrt{x}$: trasladar $3$ unidades a la derecha, escalar verticalmente por $2$, reflejar respecto del eje $x$ y trasladar $1$ unidad hacia arriba. El punto inicial es $(3,1)$, el dominio es $[3,\infty)$ y la imagen es $(-\infty,1]$.
\begin{center}
% Gráfica de la solución del ejemplo de transformaciones.
\begin{tikzpicture}[x=0.72cm,y=0.58cm,every node/.style={font=\scriptsize}]
  \draw[Gris!70] (2.3,0)--(9,0) node[right] {$x$};
  \draw[Gris!70] (3,-4.1)--(3,2) node[above] {$y$};
  \draw[Gris,dashed] (3,0)--(3,1);
  \draw[Azul,very thick,domain=3:8.5,samples=70]
    plot (\x,{-2*sqrt(\x-3)+1});
  \fill[Rojo] (3,1) circle (1.7pt);
  \node[above right] at (3,1) {$(3,1)$};
  \node[below left] at (3,0) {$3$};
  \node[left] at (3,1) {$1$};
  \node[below] at (6.5,-3.2) {$g(x)=-2\sqrt{x-3}+1$};
\end{tikzpicture}
\end{center}
\end{solucion}
\else\vspace{15mm}\fi

\section*{Operaciones y composición}

\begin{definicion}[operaciones]
Sean $f:A\to\mathbb{R}$, $g:B\to\mathbb{R}$, $\lambda\in\mathbb{R}$ y $D=A\cap B$. Se definen las siguientes operaciones:
\begin{enumerate}
  \item \textbf{Producto por un escalar:}
  \[
  \lambda f:A\to\mathbb{R},\qquad
  (\lambda f)(x)=\lambda f(x).
  \]
  \item \textbf{Suma:}
  \[
  f+g:D\to\mathbb{R},\qquad
  (f+g)(x)=f(x)+g(x).
  \]
  \item \textbf{Diferencia:}
  \[
  f-g:D\to\mathbb{R},\qquad
  (f-g)(x)=f(x)-g(x).
  \]
  \item \textbf{Producto:}
  \[
  fg:D\to\mathbb{R},\qquad
  (fg)(x)=f(x)g(x).
  \]
  \item \textbf{Cociente:} si
  \[
  D_q=\{x\in A\cap B:g(x)\neq0\},
  \]
  entonces
  \[
  \frac{f}{g}:D_q\to\mathbb{R},\qquad
  \left(\frac fg\right)(x)=\frac{f(x)}{g(x)}.
  \]
\end{enumerate}
\end{definicion}

\begin{definicion}[composición]
Sean $A,B,C\subseteq\mathbb{R}$, $g:A\to\mathbb{R}$ y $f:B\to C$. El dominio de la composición de $f$ con $g$ es
\[
D=\{x\in A:g(x)\in B\}.
\]
La función compuesta es
\[
f\circ g:D\to C,\qquad (f\circ g)(x)=f(g(x)).
\]
Para que $f\circ g$ esté definida en todo $A$, debe cumplirse
\[
\operatorname{Im}(g)\subseteq B=\operatorname{Dom}(f).
\]

\begin{center}
% Diagrama de la composición: primero actúa g y luego f.
\begin{tikzpicture}[x=1cm,y=1cm,every node/.style={font=\scriptsize}]
  \fill[Azul!5] (0,0) ellipse (0.9 and 1.15);
  \draw[Azul,thick] (0,0) ellipse (0.9 and 1.15);
  \fill[Azul!5] (3.5,0) ellipse (1.15 and 1.35);
  \draw[Azul,thick] (3.5,0) ellipse (1.15 and 1.35);
  \fill[Verde!14] (3.5,0) ellipse (0.68 and 0.88);
  \draw[Verde!75!black,thick,dashed] (3.5,0) ellipse (0.68 and 0.88);
  \fill[Azul!5] (7,0) ellipse (0.9 and 1.15);
  \draw[Azul,thick] (7,0) ellipse (0.9 and 1.15);

  \node[above] at (0,1.15) {$A=\operatorname{Dom}(g)$};
  \node[above] at (3.5,1.35) {$B=\operatorname{Dom}(f)$};
  \node at (3.5,0.62) {$\operatorname{Im}(g)$};
  \node[above] at (7,1.15) {$C$};

  \fill[Rojo] (0,0) circle (1.5pt) node[below=2pt] {$x$};
  \fill[Rojo] (3.5,-0.2) circle (1.5pt) node[below=2pt] {$g(x)$};
  \fill[Rojo] (7,0.15) circle (1.5pt) node[below=2pt] {$f(g(x))$};

  \draw[Azul,thick,->] (0.75,0.75) .. controls (1.45,1.32) and (2.15,1.42) ..
    node[above] {$g$} (2.65,0.85);
  \draw[Azul,thick,->] (4.35,0.85) .. controls (5.05,1.42) and (5.75,1.32) ..
    node[above] {$f$} (6.25,0.75);
  \draw[Verde!75!black,thick,->] (0.55,-0.92) .. controls (2.2,-2.05) and (4.8,-2.05) ..
    node[below] {$f\circ g$} (6.45,-0.92);
  \draw[Rojo,thin,->] (0.12,0.02) .. controls (1.2,-0.48) and (2.25,-0.48) .. (3.38,-0.22);
  \draw[Rojo,thin,->] (3.62,-0.18) .. controls (4.75,-0.55) and (5.85,-0.45) .. (6.88,0.12);
\end{tikzpicture}
\end{center}
El óvalo verde representa $\operatorname{Im}(g)$. Al estar contenido en $\operatorname{Dom}(f)$, cada salida de $g$ puede utilizarse como entrada de $f$.
\end{definicion}

En general, $g\circ f\neq f\circ g$: puede cambiar la regla, el dominio o ambos.

\begin{ejemplo}[composición y dominio]
Sean $f:\mathbb{R}\to\mathbb{R}$, $f(x)=x^2+1$, y $g:[0,\infty)\to\mathbb{R}$, $g(x)=\sqrt{x}$. Calcule $g\circ f$ y $f\circ g$ e indique sus dominios.
\end{ejemplo}
\ifmostrarSoluciones
\begin{solucion}
Como $x^2+1\geq1$, se tiene
\[
(g\circ f)(x)=\sqrt{x^2+1},\qquad
\operatorname{Dom}(g\circ f)=\mathbb{R}.
\]
En el otro orden, la función interior exige $x\geq0$:
\[
(f\circ g)(x)=(\sqrt{x})^2+1=x+1,\qquad
\operatorname{Dom}(f\circ g)=[0,\infty).
\]
La simplificación de $(\sqrt{x})^2$ no elimina la restricción original $x\geq0$.
\end{solucion}
\else\vspace{15mm}\fi

\section*{Simetría y monotonía de una función}

\begin{definicion}[conjunto simétrico]
Un conjunto $A\subseteq\mathbb{R}$ es \textbf{simétrico respecto de cero} si
\[
\forall x\in A,\qquad -x\in A.
\]
Esta condición garantiza que, junto con cada entrada $x$, también pueda evaluarse la entrada opuesta $-x$.
\end{definicion}

\begin{definicion}[función par e impar]
Sea $f:A\to\mathbb{R}$, con $A$ simétrico respecto de cero.
\begin{itemize}
  \item $f$ es \textbf{par} si, para todo $x\in A$, $f(-x)=f(x)$. Su gráfico es simétrico respecto del eje $y$.
  \item $f$ es \textbf{impar} si, para todo $x\in A$, $f(-x)=-f(x)$. Su gráfico es simétrico respecto del origen.
\end{itemize}
\end{definicion}

\begin{ejemplo}[gráficas de funciones par e impar]
La función $p(x)=x^2$ es par: los puntos correspondientes a $x$ y $-x$ tienen la misma ordenada. La función $q(x)=x^3$ es impar: los puntos correspondientes son simétricos respecto del origen.
\begin{center}
\setlength{\tabcolsep}{3mm}
\begin{tabular}{cc}
% Ejemplo de función par: simetría respecto del eje y.
\begin{tikzpicture}[x=0.75cm,y=0.48cm,every node/.style={font=\scriptsize}]
  \draw[Gris!70] (-1.9,0)--(1.9,0) node[right] {$x$};
  \draw[Gris!70] (0,-0.25)--(0,2.8) node[above] {$y$};
  \draw[Azul,very thick,domain=-1.6:1.6,samples=60] plot (\x,{\x*\x});
  \draw[Rojo,dashed] (-1,1)--(1,1);
  \draw[Rojo,dashed] (-1,0)--(-1,1) (1,0)--(1,1);
  \fill[Rojo] (-1,0) circle (1.2pt) (1,0) circle (1.2pt);
  \fill[Rojo] (-1,1) circle (1.5pt) (1,1) circle (1.5pt);
  \node[below] at (-1,0) {$-x$};
  \node[below] at (1,0) {$x$};
  \node[below=3pt] at (0,-0.72) {$p(x)=x^2$ (par)};
\end{tikzpicture}
&
% Ejemplo de función impar: simetría respecto del origen.
\begin{tikzpicture}[x=0.75cm,y=0.48cm,every node/.style={font=\scriptsize}]
  \draw[Gris!70] (-1.9,0)--(1.9,0) node[right] {$x$};
  \draw[Gris!70] (0,-2.45)--(0,2.45) node[above] {$y$};
  \draw[Azul,very thick,domain=-1.3:1.3,samples=60] plot (\x,{\x*\x*\x});
  \draw[Rojo,dashed] (-1,-1)--(1,1);
  \draw[Rojo,dashed] (-1,0)--(-1,-1) (1,0)--(1,1);
  \fill[Rojo] (-1,0) circle (1.2pt) (1,0) circle (1.2pt);
  \fill[Rojo] (-1,-1) circle (1.5pt) (1,1) circle (1.5pt);
  \node[above] at (-1,0) {$-x$};
  \node[below] at (1,0) {$x$};
  \node[below=3pt] at (0,-2.45) {$q(x)=x^3$ (impar)};
\end{tikzpicture}
\end{tabular}
\end{center}
En efecto, $p(-x)=(-x)^2=p(x)$, mientras que $q(-x)=(-x)^3=-q(x)$.
\end{ejemplo}

\begin{definicion}[monotonía de una función]
Sea $f:A\to\mathbb{R}$ e $I\subseteq A$ un intervalo. Para todo $u,v\in I$:
\begin{itemize}
  \item $f$ es \textbf{creciente} en $I$ si $u\leq v$ implica $f(u)\leq f(v)$;
  \item $f$ es \textbf{estrictamente creciente} en $I$ si $u<v$ implica $f(u)<f(v)$;
  \item $f$ es \textbf{decreciente} en $I$ si $u\leq v$ implica $f(u)\geq f(v)$;
  \item $f$ es \textbf{estrictamente decreciente} en $I$ si $u<v$ implica $f(u)>f(v)$.
\end{itemize}
Se dice que $f$ es \textbf{monótona} en $I$ si es creciente o decreciente en ese intervalo.

\begin{center}
\setlength{\tabcolsep}{3mm}
\renewcommand{\arraystretch}{1.25}
\begin{tabular}{cc}
% Función creciente, con un tramo constante.
\begin{tikzpicture}[x=0.78cm,y=0.52cm,every node/.style={font=\scriptsize}]
  \draw[Gris!70] (-1.9,0)--(1.9,0) node[right] {$x$};
  \draw[Gris!70] (0,-1.25)--(0,1.45) node[above] {$y$};
  \draw[Azul,very thick] (-1.65,-1)--(-0.8,0.15)--(0.8,0.15)--(1.65,1.2);
  \draw[Rojo,dashed] (-1.2,0)--(-1.2,-0.39) (1.2,0)--(1.2,0.64);
  \fill[Rojo] (-1.2,-0.39) circle (1.4pt) (1.2,0.64) circle (1.4pt);
  \node[above] at (-1.2,0) {$u$};
  \node[below] at (1.2,0) {$v$};
  \node[left] at (-1.2,-0.39) {$f(u)$};
  \node[right] at (1.2,0.64) {$f(v)$};
  \node[below=2pt,align=center] at (0,-1.25) {creciente\\$u<v\Rightarrow f(u)\leq f(v)$};
\end{tikzpicture}
&
% Función estrictamente creciente.
\begin{tikzpicture}[x=0.78cm,y=0.52cm,every node/.style={font=\scriptsize}]
  \draw[Gris!70] (-1.9,0)--(1.9,0) node[right] {$x$};
  \draw[Gris!70] (0,-1.25)--(0,1.45) node[above] {$y$};
  \draw[Azul,very thick] (-1.65,-1.05)--(1.65,1.2);
  \draw[Rojo,dashed] (-0.8,0)--(-0.8,-0.47) (0.8,0)--(0.8,0.62);
  \fill[Rojo] (-0.8,-0.47) circle (1.4pt) (0.8,0.62) circle (1.4pt);
  \node[above] at (-0.8,0) {$u$};
  \node[below] at (0.8,0) {$v$};
  \node[left] at (-0.8,-0.47) {$f(u)$};
  \node[right] at (0.8,0.62) {$f(v)$};
  \node[below=2pt,align=center] at (0,-1.25) {estrictamente creciente\\$u<v\Rightarrow f(u)<f(v)$};
\end{tikzpicture}
\\
% Función decreciente, con un tramo constante.
\begin{tikzpicture}[x=0.78cm,y=0.52cm,every node/.style={font=\scriptsize}]
  \draw[Gris!70] (-1.9,0)--(1.9,0) node[right] {$x$};
  \draw[Gris!70] (0,-1.25)--(0,1.45) node[above] {$y$};
  \draw[Azul,very thick] (-1.65,1.2)--(-0.8,0.15)--(0.8,0.15)--(1.65,-1);
  \draw[Rojo,dashed] (-1.2,0)--(-1.2,0.64) (1.2,0)--(1.2,-0.39);
  \fill[Rojo] (-1.2,0.64) circle (1.4pt) (1.2,-0.39) circle (1.4pt);
  \node[below] at (-1.2,0) {$u$};
  \node[above] at (1.2,0) {$v$};
  \node[left] at (-1.2,0.64) {$f(u)$};
  \node[right] at (1.2,-0.39) {$f(v)$};
  \node[below=2pt,align=center] at (0,-1.25) {decreciente\\$u<v\Rightarrow f(u)\geq f(v)$};
\end{tikzpicture}
&
% Función estrictamente decreciente.
\begin{tikzpicture}[x=0.78cm,y=0.52cm,every node/.style={font=\scriptsize}]
  \draw[Gris!70] (-1.9,0)--(1.9,0) node[right] {$x$};
  \draw[Gris!70] (0,-1.25)--(0,1.45) node[above] {$y$};
  \draw[Azul,very thick] (-1.65,1.2)--(1.65,-1.05);
  \draw[Rojo,dashed] (-0.8,0)--(-0.8,0.62) (0.8,0)--(0.8,-0.47);
  \fill[Rojo] (-0.8,0.62) circle (1.4pt) (0.8,-0.47) circle (1.4pt);
  \node[below] at (-0.8,0) {$u$};
  \node[above] at (0.8,0) {$v$};
  \node[left] at (-0.8,0.62) {$f(u)$};
  \node[right] at (0.8,-0.47) {$f(v)$};
  \node[below=2pt,align=center] at (0,-1.25) {estrictamente decreciente\\$u<v\Rightarrow f(u)>f(v)$};
\end{tikzpicture}
\end{tabular}
\end{center}
Los tramos horizontales son compatibles con la monotonía no estricta, pero no con la monotonía estricta.
\end{definicion}

\begin{nota}
Una función puede cambiar de comportamiento entre intervalos. Por ello, afirmar que es creciente o decreciente siempre debe acompañarse del intervalo donde se verifica la propiedad.
\end{nota}

\begin{ejemplo}[paridad y monotonía]
Analice la paridad y la monotonía de $r(x)=|x|$.
\end{ejemplo}
\ifmostrarSoluciones
\begin{solucion}
El dominio $\mathbb{R}$ es simétrico y
\[
r(-x)=|-x|=|x|=r(x),
\]
por lo que $r$ es par. Además, es estrictamente decreciente en $(-\infty,0]$ y estrictamente creciente en $[0,\infty)$; no es monótona en todo $\mathbb{R}$.
\end{solucion}
\else\vspace{15mm}\fi

\begin{ejercicios}[transformaciones, operaciones y comportamiento]
\begin{enumerate}
  \item Describa cómo obtener a partir de $f(x)=x^2$ las gráficas de $f_1(x)=(x-3)^2+2$ y $f_2(x)=-\tfrac12(x+1)^2+4$.
  \item Determine el mayor conjunto $D\subseteq\mathbb{R}$ para que $g(x)=3\sqrt{2-x}-1$ defina una función real $g:D\to\mathbb{R}$; luego identifique sus transformaciones y su imagen.
  \item Bosqueje $f(x)=|x+2|-3$ y halle sus intersecciones con los ejes.
  %\item Si el gráfico de $f$ contiene $(2,-1)$, determine puntos correspondientes en $f(x-3)+4$, $-f(x)$ y $f(-x)$.
  \item Sean $f:\mathbb{R}\to\mathbb{R}$, $f(x)=x^2-1$, y $g:\mathbb{R}\to\mathbb{R}$, $g(x)=x+2$. Calcule $3f$, $f+g$, $fg$ y $f/g$, e indique el dominio de cada función resultante.
  \item Sean $f:[1,\infty)\to\mathbb{R}$, $f(x)=\sqrt{x-1}$, y $g:\mathbb{R}\to\mathbb{R}$, $g(x)=x^2$. Halle ambas composiciones e indique sus dominios.
  \item Descomponga como $g\circ f$:
  \[
  \sqrt{4-x^2},\qquad \frac{1}{(2x+3)^2},\qquad |x^3-1|.
  \]
  \item Clasifique como par, impar o ninguna: $x^6-2x^2$, $x^5-3x$, $x^2+x$ y $|x|+2$.
  \item Indique intervalos de crecimiento y decrecimiento para $q(x)=|x-2|$.
  \item Sea $f:\mathbb{R}\to\mathbb{R}$, $f(x)=2|x+1|-3$. Analice su imagen, sus intersecciones con los ejes, su monotonía y su simetría.
\end{enumerate}
\end{ejercicios}

\ifmostrarSoluciones
\begin{solucion}[de los ejercicios]
\begin{enumerate}
  \item $(x-3)^2+2$: trasladar $3$ a la derecha y $2$ arriba. Para $-\tfrac12(x+1)^2+4$: trasladar $1$ a la izquierda, comprimir verticalmente por $1/2$, reflejar en el eje $x$ y trasladar $4$ arriba.
  \begin{center}
  \setlength{\tabcolsep}{2mm}
  \begin{tabular}{cc}
  % Gráfica de la primera función del ejercicio 1.
  \begin{tikzpicture}[x=0.55cm,y=0.34cm,every node/.style={font=\scriptsize}]
    \draw[Gris!70] (0.3,0)--(5.8,0) node[right] {$x$};
    \draw[Gris!70] (0.5,-0.3)--(0.5,7.1) node[above] {$y$};
    \draw[Gris,dashed] (3,0)--(3,6.8);
    \draw[Azul,very thick,domain=0.8:5.2,samples=60] plot (\x,{(\x-3)*(\x-3)+2});
    \fill[Rojo] (3,2) circle (1.6pt);
    \node[right] at (3,2) {$(3,2)$};
    \node[below] at (3,0) {$3$};
    \node[below=4pt,align=center] at (3,-0.3) {$f_1(x)=(x-3)^2+2$};
  \end{tikzpicture}
  &
  % Gráfica de la segunda función del ejercicio 1.
  \begin{tikzpicture}[x=0.55cm,y=0.42cm,every node/.style={font=\scriptsize}]
    \draw[Gris!70] (-4.4,0)--(2.4,0) node[right] {$x$};
    \draw[Gris!70] (0,-0.8)--(0,4.8) node[above] {$y$};
    \draw[Gris,dashed] (-1,-0.6)--(-1,4.6);
    \draw[Azul,very thick,domain=-4.1:2.1,samples=60] plot (\x,{-0.5*(\x+1)*(\x+1)+4});
    \fill[Rojo] (-1,4) circle (1.6pt);
    \node[right] at (-1,4) {$(-1,4)$};
    \node[below] at (-1,0) {$-1$};
    \node[below=4pt,align=center] at (-1,-0.8) {$f_2(x)=-\tfrac12(x+1)^2+4$};
  \end{tikzpicture}
  \end{tabular}
  \end{center}
  \item Dominio $(-\infty,2]$ e imagen $[-1,\infty)$.
  \item Vértice $(-2,-3)$; ceros $x=-5,1$; intersección vertical $(0,-1)$.
  %\item Los puntos son $(5,3)$, $(2,1)$ y $(-2,-1)$, respectivamente.
  \item
  \[
  3f=3x^2-3,\qquad f+g=x^2+x+1,
  \]
  \[
  fg=x^3+2x^2-x-2,
  \]
  todas con dominio $\mathbb{R}$, y $f/g=(x^2-1)/(x+2)$ con dominio $\mathbb{R}\setminus\{-2\}$.
  \item
  \[
  (f\circ g)(x)=\sqrt{x^2-1},\quad \operatorname{Dom}=(-\infty,-1]\cup[1,\infty),
  \]
  \[
  (g\circ f)(x)=x-1,\quad \operatorname{Dom}=[1,\infty).
  \]
  \item Use $f_1(x)=4-x^2$, $g_1(u)=\sqrt{u}$; $f_2(x)=2x+3$, $g_2(u)=1/u^2$; y $f_3(x)=x^3-1$, $g_3(u)=|u|$.
  \item Par, impar, ninguna y par, respectivamente.
  \item Decreciente en $(-\infty,2]$ y creciente en $[2,\infty)$.
  \item Dominio $\mathbb{R}$, imagen $[-3,\infty)$, ceros $-5/2$ y $1/2$, e intersección vertical $(0,-1)$. Es decreciente hasta $-1$, creciente desde $-1$ y simétrica respecto de $x=-1$; no es par ni impar.
\end{enumerate}
\end{solucion}
\fi

\ifdefined\documentoSemanal
\else
\fuentes{\textbf{Para ampliar:} \emph{Cálculo en una variable: notas de clase 2018-B}, Resumen No. 1, pp. 2--4; J. Stewart, \emph{Cálculo de una variable: trascendentes tempranas}, §1.3; y G. Strang--E. Herman, \emph{Cálculo, volumen 1}, OpenStax, §§1.1--1.2.}
\end{document}
\fi
